% =====================================================================
%  PSH-PR -- Manual Operativo do Projeto (MOP)
%  Modelo de estrutura metodológica -- Subcomponente 3.2.b
%  Instituto Água e Terra (IAT)
%
%  Compilar com:  latexmk -pdf main.tex
% =====================================================================
\documentclass[
  12pt,
  oneside,
  a4paper,
  english,
  brazil
]{abntex2}

% Mostrar (true) ou ocultar (false) as notas de pendência no PDF.
% Para a versão final, trocar \mostrarpendenciastrue por \mostrarpendenciasfalse.
\newif\ifmostrarpendencias
\mostrarpendenciastrue

% =====================================================================
%  Pacotes e configurações gerais
% =====================================================================

% --- Fontes e codificação ---------------------------------------------
\usepackage{lmodern}
\usepackage[T1]{fontenc}
\usepackage[utf8]{inputenc}
\usepackage{microtype}

% --- Parágrafos, listas e tabelas -------------------------------------
\usepackage{indentfirst}
\usepackage[table]{xcolor}
\usepackage{graphicx}
\usepackage{array}
\usepackage{tabularx}
\usepackage{booktabs}
\usepackage{multirow}
\usepackage{pdflscape}
\usepackage{enumitem}
\setlist{nosep, leftmargin=*, itemsep=2pt}

% --- Figuras (fluxogramas e cronograma) -------------------------------
\usepackage{tikz}
\usetikzlibrary{positioning, arrows.meta, shapes.geometric, shapes.symbols, calc, fit, backgrounds}
\usepackage{pgfgantt}
\usepackage{adjustbox}

% --- Notas de pendência (revisão do MOP) ------------------------------
\ifmostrarpendencias
  \usepackage[colorinlistoftodos, textsize=footnotesize,
              figwidth=\linewidth]{todonotes}
\else
  \usepackage[disable]{todonotes}
\fi

% --- Citações (ABNT, autor-data) --------------------------------------
\usepackage[alf, abnt-emphasize=bf]{abntex2cite}

% --- Cores do documento -----------------------------------------------
\definecolor{azulpsh}{RGB}{0,84,147}
\definecolor{azulclaro}{RGB}{222,235,247}
\definecolor{verdeclaro}{RGB}{226,239,218}
\definecolor{laranjaclaro}{RGB}{252,228,214}
\definecolor{cinzaclaro}{RGB}{242,242,242}
% fases do cronograma do POA (Anexo A)
\definecolor{corTR}{RGB}{189,215,238}
\definecolor{corLicitacao}{RGB}{255,192,0}
\definecolor{corExecucao}{RGB}{0,84,147}
\definecolor{corInterna}{RGB}{112,173,71}

% --- Hiperlinks -------------------------------------------------------
\makeatletter
\hypersetup{
  pdftitle={\@title},
  pdfauthor={\@author},
  pdfsubject={Subcomponente 3.2.b -- estrutura metodológica},
  pdfkeywords={PSH-PR}{MOP}{saneamento rural}{IAT},
  colorlinks=true,
  linkcolor=azulpsh,
  citecolor=azulpsh,
  filecolor=azulpsh,
  urlcolor=azulpsh,
  bookmarksdepth=4
}
\makeatother

% --- Quadros (ABNT NBR 14724) -----------------------------------------
% A ABNT distingue "quadro" (conteúdo textual) de "tabela" (dados
% numéricos). O abnTeX2 não traz o ambiente pronto; ele é criado aqui.
\newcommand{\quadroname}{Quadro}
\newcommand{\listofquadrosname}{Lista de quadros}
\newfloat[chapter]{quadro}{loq}{\quadroname}
\newlistof{listofquadros}{loq}{\listofquadrosname}
\newlistentry{quadro}{loq}{0}
\setfloatadjustment{quadro}{\centering}
\counterwithout{quadro}{chapter}
\renewcommand{\cftquadroname}{\quadroname\space}
\renewcommand*{\cftquadroaftersnum}{\hfill--\hfill}
\setfloatlocations{quadro}{htbp}

% Numeração contínua de figuras e tabelas (sem prefixo do capítulo)
\counterwithout{figure}{chapter}
\counterwithout{table}{chapter}

% Numeração e sumário até \subsubsection (as fases, em \paragraph, não
% são numeradas nem entram no sumário)
\setsecnumdepth{subsubsection}
\settocdepth{subsubsection}

% --- Espaçamentos -----------------------------------------------------
\setlength{\parindent}{1.25cm}
\setlength{\parskip}{0.2cm}

% =====================================================================
%  Comandos e ambientes próprios do MOP
%  (usados em todas as atividades -- manter a mesma forma de uso)
% =====================================================================

% ---------------------------------------------------------------------
% \pendencia{texto}
%   Registra, no ponto do texto, uma inconsistência ou lacuna a
%   resolver. Todas aparecem listadas no Apêndice A. Para ocultá-las
%   na versão final, usar \mostrarpendenciasfalse em main.tex.
% ---------------------------------------------------------------------
\newcommand{\pendencia}[1]{%
  \todo[inline, color=laranjaclaro, bordercolor=orange,
        caption={#1}]{\textbf{Pendência:} #1}}

% Lista das pendências (sem título próprio), usada no Apêndice A
\makeatletter
\newcommand{\listapendencias}{\@starttoc{tdo}}
\makeatother

% ---------------------------------------------------------------------
% \fluxograma{arquivo}
%   Insere um fluxograma da pasta figuras/, reduzindo-o apenas se for
%   mais largo que a mancha de texto.
% ---------------------------------------------------------------------
\newcommand{\fluxograma}[1]{\adjustbox{max width=\linewidth}{\input{figuras/#1}}}

% ---------------------------------------------------------------------
% \atividade{3.2.2.x}
%   Grafia padronizada do código de uma atividade no texto corrido.
% ---------------------------------------------------------------------
\newcommand{\atividade}[1]{Atividade~#1}

% ---------------------------------------------------------------------
% Cabeçalho de quadros/tabelas (linha colorida em negrito)
% ---------------------------------------------------------------------
\newcommand{\cab}[1]{\cellcolor{azulclaro}\textbf{#1}}

% Colunas utilitárias para tabularx
\newcolumntype{L}{>{\raggedright\arraybackslash}X}
\newcolumntype{C}{>{\centering\arraybackslash}X}
\newcolumntype{P}[1]{>{\raggedright\arraybackslash}p{#1}}
\newcolumntype{Q}[1]{>{\centering\arraybackslash}p{#1}}

% ---------------------------------------------------------------------
% Estilos TikZ dos fluxogramas
% ---------------------------------------------------------------------
\tikzset{
  etapa/.style={draw=azulpsh, fill=azulclaro, rounded corners=2pt,
                align=center, font=\footnotesize, inner sep=4pt,
                minimum height=0.9cm, text width=#1},
  etapa/.default={3.4cm},
  marco/.style={draw=azulpsh, fill=azulpsh, text=white, rounded corners=2pt,
                align=center, font=\footnotesize\bfseries, inner sep=4pt,
                minimum height=0.9cm, text width=#1},
  marco/.default={7.4cm},
  externo/.style={draw=gray, fill=cinzaclaro, rounded corners=2pt,
                  align=center, font=\footnotesize, inner sep=4pt,
                  minimum height=0.9cm, text width=#1},
  externo/.default={3.4cm},
  decisao/.style={draw=orange!80!black, fill=laranjaclaro, diamond, aspect=2.2,
                  align=center, font=\footnotesize, inner sep=1pt},
  correcao/.style={draw=orange!80!black, fill=laranjaclaro, rounded corners=2pt,
                   align=center, font=\footnotesize, inner sep=4pt,
                   minimum height=0.9cm, text width=#1},
  correcao/.default={2.6cm},
  transversal/.style={draw=verdeclaro!50!black, fill=verdeclaro, dashed,
                      rounded corners=2pt, align=center, font=\footnotesize,
                      inner sep=4pt, text width=#1},
  seta/.style={-{Stealth[length=2.2mm]}, thick, draw=azulpsh!80!black},
  setatracejada/.style={-{Stealth[length=2.2mm]}, thick, dashed, draw=gray},
  rotulo/.style={font=\scriptsize\itshape, fill=white, inner sep=1pt},
  % etapa em forma de seta (linhas do tempo)
  fase/.style={signal, signal to=east, signal from=west, draw=azulpsh,
               fill=azulclaro, align=center, font=\scriptsize\bfseries,
               minimum height=1.1cm, text width=#1, inner sep=2pt},
  fase/.default={2.1cm},
  % caixa de produto ou documento resultante de uma etapa
  produto/.style={draw=azulpsh!70, fill=white, align=center, font=\scriptsize,
                  text width=#1, inner sep=3pt, minimum height=0.7cm,
                  rounded corners=1pt},
  produto/.default={2.1cm},
  % faixa (raia) de responsável em fluxogramas do tipo swimlane
  raia/.style={draw=gray!60, fill=cinzaclaro!70, minimum width=#1,
               minimum height=0.9cm, anchor=north west},
  titulo raia/.style={font=\scriptsize\bfseries, text=azulpsh, align=center,
                      text width=1.6cm}
}


% ---------------------------------------------------------------------
% Dados do documento
% ---------------------------------------------------------------------
\titulo{Manual Operativo do Projeto de Segurança Hídrica do Paraná (PSH-PR)}
\autor{Instituto Água e Terra -- IAT}
\local{Curitiba}
\data{2026}
\instituicao{Governo do Estado do Paraná\par Instituto Água e Terra}
\tipotrabalho{Manual Operativo}
\preambulo{Subcomponente 3.2.b -- Implantação e reabilitação de sistemas
coletivos de abastecimento de água e esgotamento sanitário no meio rural.
Modelo de estrutura metodológica.}

\begin{document}

\selectlanguage{brazil}
\frenchspacing

% ---------------------------------------------------------------------
% Elementos pré-textuais
% ---------------------------------------------------------------------
\pretextual
\imprimirfolhaderosto

\pdfbookmark[0]{\listfigurename}{lof}
\listoffigures*
\cleardoublepage

\pdfbookmark[0]{\listofquadrosname}{loq}
\listofquadros*
\cleardoublepage

\pdfbookmark[0]{\listtablename}{lot}
\listoftables*
\cleardoublepage

\begin{siglas}
  \item[BIRD] Banco Internacional para Reconstrução e Desenvolvimento (Banco Mundial)
  \item[ESMF] \textit{Environmental and Social Management Framework}
  \item[IAT] Instituto Água e Terra
  \item[IDR-Paraná] Instituto de Desenvolvimento Rural do Paraná
  \item[IPDM] Índice Ipardes de Desempenho Municipal
  \item[LMP] \textit{Labor Management Procedures}
  \item[MGAS] Marco de Gestão Ambiental e Social
  \item[MQR] Mecanismo de Queixas e Reclamações
  \item[MRF] Marco de Política de Reassentamento
  \item[O\&M] Operação e manutenção
  \item[PEPI] Plano de Engajamento das Partes Interessadas
  \item[PGR] Procedimentos de Gestão de Mão de Obra
  \item[PIP] Plano Integrado de Propriedade
  \item[PSH-PR] Projeto de Segurança Hídrica do Paraná
  \item[RWSS] \textit{Rural Water Supply and Sanitation}
  \item[SECID] Secretaria de Estado das Cidades
  \item[SEP] \textit{Stakeholder Engagement Plan}
  \item[TR] Termo de Referência
  \item[TRD] Termo de Recebimento Definitivo
  \item[UGP] Unidade de Gerenciamento do Programa
\end{siglas}

\pdfbookmark[0]{\contentsname}{toc}
\tableofcontents*
\cleardoublepage

% ---------------------------------------------------------------------
% Elementos textuais
% ---------------------------------------------------------------------
\textual

% =====================================================================
%  Apresentação do modelo de estrutura
% =====================================================================
\chapter*[Apresentação do modelo]{Apresentação do modelo}
\addcontentsline{toc}{chapter}{Apresentação do modelo}

Este documento propõe uma nova forma de apresentar as ações do Instituto
Água e Terra (IAT) no Manual Operativo do Projeto (MOP) do PSH-PR. Toma-se
como piloto o Subcomponente 3.2.b -- \textit{Implantação e reabilitação de
sistemas coletivos de abastecimento de água e esgotamento sanitário no meio
rural} (item B.3.2.2 do MOP de 17 de março de 2026). O conteúdo técnico
foi preservado; o que muda é a organização. As atividades passam a ser
descritas como o capítulo de metodologia de uma dissertação: justificativa
e objetivos, área de estudo e critérios de seleção, delineamento, procedimentos
por atividade, salvaguardas, monitoramento e cronograma.

A mesma estrutura deve ser replicada nos demais subcomponentes sob
responsabilidade do IAT (Componente~1 e Subcomponente~3.2.b). O arquivo
\texttt{modelos/modelo-atividade.tex} contém o esqueleto em branco de uma
atividade.

\section*{Estrutura padrão do capítulo}

\begin{enumerate}[label=\arabic*.]
  \item \textbf{Contextualização e objetivos}: justificativa, objetivo geral,
        objetivos específicos e metas físicas.
  \item \textbf{Área de abrangência e critérios de seleção}: onde a ação
        acontece e como são escolhidos os territórios e beneficiários.
  \item \textbf{Delineamento metodológico}: premissas, encadeamento das
        atividades (fluxograma) e arranjo institucional e financeiro.
  \item \textbf{Procedimentos metodológicos por atividade}: uma seção por
        atividade, sempre com as mesmas subdivisões:
        \begin{enumerate}[label=\alph*)]
          \item objetivo e fundamentação;
          \item ficha-síntese (quadro padronizado);
          \item procedimentos, organizados em fases;
          \item etapas, produtos e evidências (quadro detalhado, quando houver);
          \item interfaces com outras atividades.
        \end{enumerate}
  \item \textbf{Salvaguardas ambientais e sociais}: riscos, mitigação e
        instrumentos do Quadro Ambiental e Social do Banco Mundial.
  \item \textbf{Monitoramento, avaliação e indicadores}: indicadores, metas,
        fontes de verificação e fluxo de validação.
  \item \textbf{Cronograma de execução}: diagrama de Gantt.
\end{enumerate}

\section*{Convenções de redação e formatação}

\begin{itemize}
  \item Redação impessoal, no futuro do presente (\textit{serão realizadas},
        \textit{será elaborado}), como em um projeto de pesquisa.
  \item Informação textual vai em \textbf{quadro} e dados numéricos vão em
        \textbf{tabela} (ABNT NBR 14724). Todo quadro, tabela e figura traz
        a fonte.
  \item Os códigos das atividades (3.2.2.1 a 3.2.2.6) seguem o MOP e são
        citados com o comando \verb|\atividade{3.2.2.x}|.
  \item Inconsistências encontradas durante a reestruturação são marcadas
        com \verb|\pendencia{...}| e listadas no Apêndice~\ref{ap:pendencias}.
        Para ocultá-las na versão final, basta desativar a opção em
        \texttt{main.tex}.
\end{itemize}

\section*{Correspondência com o MOP vigente}

O \autoref{qua:correspondencia} indica onde o conteúdo do item B.3.2.2 do MOP
foi reposicionado na nova estrutura.

\begin{quadro}[htbp]
\caption{Correspondência entre o MOP (versão de 17/03/2026) e a estrutura proposta}
\label{qua:correspondencia}
\small
\renewcommand{\arraystretch}{1.2}
\begin{tabularx}{\textwidth}{|L|L|}
\hline
\cab{MOP -- item B.3.2.2} & \cab{Estrutura proposta} \\ \hline
Texto introdutório e objetivos específicos & \ref{sec:contexto} Contextualização e objetivos \\ \hline
Frentes de trabalho (83 comunidades; 2.300 sistemas) & \ref{subsec:metas} Metas \\ \hline
Área de abrangência e critérios de seleção & \ref{sec:area} Área de abrangência e critérios de seleção \\ \hline
Executor, forma de execução, prazo e custo & \ref{subsec:arranjo} Arranjo institucional e financiamento \\ \hline
Quadro 26 -- Aspectos socioambientais & \ref{sec:salvaguardas} Salvaguardas ambientais e sociais \\ \hline
Quadro 27 -- Indicadores & \ref{sec:monitoramento} Monitoramento, avaliação e indicadores \\ \hline
B.3.2.2.1 a B.3.2.2.6 (atividades) & \ref{sec:procedimentos} Procedimentos metodológicos por atividade \\ \hline
Quadros 28 e 29 -- Etapas das atividades & Quadros de etapas dentro de cada atividade \\ \hline
Quadro 30 -- Atividades, responsáveis e prazos & Fichas-síntese de cada atividade e \ref{sec:cronograma} Cronograma \\ \hline
Quadro 31 -- Produtos e evidências & Fichas-síntese de cada atividade \\ \hline
\end{tabularx}
\fonte{Elaborado pelo IAT (2026).}
\end{quadro}

% =====================================================================
%  Subcomponente 3.2.b
%  Implantação e reabilitação de sistemas coletivos de abastecimento de
%  água e esgotamento sanitário no meio rural  (MOP, item B.3.2.2)
% =====================================================================
\chapter[Subcomponente 3.2.b -- Abastecimento e esgotamento no meio rural]%
{Subcomponente 3.2.b -- Implantação e reabilitação de sistemas coletivos de
abastecimento de água e esgotamento sanitário no meio rural}
\label{cap:3-2-b}

% =====================================================================
%  1. Contextualização e objetivos
% =====================================================================
\section{Contextualização e objetivos}
\label{sec:contexto}

\subsection{Justificativa}

O Subcomponente 3.2 do PSH-PR tem como objetivo ampliar o acesso a serviços
sustentáveis de abastecimento de água e esgotamento sanitário no meio rural
paranaense, contribuindo para a universalização do saneamento, a promoção da
saúde pública e a proteção dos recursos hídricos \cite{mop2026}. A
universalização, estabelecida pelo novo marco legal do saneamento básico
\cite{lei14026}, ainda está distante nas áreas rurais do Estado: segundo dados
do \citeonline{sinisa2024}, a cobertura por rede de abastecimento de água
alcança 99,47\% da população urbana, mas apenas 15,11\% da população rural; no
esgotamento sanitário, a cobertura por rede coletora é de 84,41\% na área
urbana e de 2,04\% na área rural.

O saneamento rural parte do propósito de universalizar o acesso aos serviços
por meio de ações que garantam equidade, integralidade, intersetorialidade,
sustentabilidade dos serviços, participação e controle social. Diferentemente
do meio urbano, a alta dispersão populacional exige soluções técnicas --
coletivas ou individuais -- adaptadas à realidade local e acompanhadas de
modelos de gestão capazes de manter os sistemas em funcionamento após a sua
implantação \cite{mop2026}.

\pendencia{No MOP, o primeiro parágrafo do item B.3.2.2 menciona apenas
``universalizar o acesso ao Esgotamento Sanitário'', embora a ação trate
também de abastecimento de água. Há ainda um erro de digitação:
``soluções técnicas 0(coletivas ou individuais)''.}

\subsection{Objetivo geral}

Ampliar o acesso da população rural do Paraná ao abastecimento de água e ao
esgotamento sanitário por meio da implantação e da reabilitação de sistemas
coletivos de abastecimento de água, da instalação de sistemas descentralizados
de esgotamento sanitário e da implantação de modelos de gestão que assegurem
a sustentabilidade dos serviços.

\subsection{Objetivos específicos}

\begin{enumerate}[label=\alph*)]
  \item desenvolver modelos de gestão que permitam alcançar a
        universalização do saneamento rural com qualidade e sustentabilidade
        na prestação dos serviços públicos;
  \item implantar novos sistemas coletivos de abastecimento de água e
        sistemas descentralizados de esgotamento sanitário em comunidades
        rurais do Estado, visando ao princípio da integralidade;
  \item instalar sistemas descentralizados de esgotamento sanitário nas
        áreas prioritárias do PSH-PR de atuação do IDR-Paraná,
        preferencialmente nas áreas de mananciais de abastecimento público;
  \item reabilitar sistemas coletivos de abastecimento de água e substituir
        sistemas descentralizados de esgotamento sanitário em comunidades
        rurais, com foco na recuperação e na sustentabilidade.
\end{enumerate}

\subsection{Metas}
\label{subsec:metas}

Após a aplicação dos critérios de seleção definidos em âmbito estadual
(\autoref{sec:area}), a ação será conduzida em três frentes de trabalho
\cite{mop2026}:

\begin{enumerate}[label=\roman*)]
  \item implantar ou reabilitar sistemas de abastecimento de água em
        83 comunidades rurais (5.810 domicílios) e instalar sistemas
        descentralizados de esgotamento sanitário em 70\% dos domicílios
        atendidos (4.067 domicílios);
  \item aplicar modelos de gestão aos sistemas de abastecimento de água e de
        esgotamento sanitário implantados nas comunidades rurais e aos
        sistemas de abastecimento implantados pelo Programa Sanepar Rural;
  \item instalar 2.300 sistemas descentralizados de esgotamento sanitário,
        preferencialmente em mananciais de abastecimento público
        selecionados por critério qualitativo, na área prioritária do PSH-PR
        de atuação do IDR-Paraná.
\end{enumerate}

No conjunto, estima-se atender 22.708 habitantes com acesso à água potável e
ao esgotamento sanitário (\autoref{qua:indicadores}).

\pendencia{O MOP associa os 22.708 habitantes às 83 comunidades. Pelo fator
de 2,8 habitantes por domicílio (Censo 2022), adotado nos indicadores do
Componente~3, as 83 comunidades (5.810 domicílios) somam cerca de 16.268
habitantes. O total de 22.708 corresponde a (5.810 + 2.300) domicílios
$\times$ 2,8, ou seja, inclui os domicílios da área do IDR-Paraná. Ajustar o
texto do MOP.}

% =====================================================================
%  2. Área de abrangência e critérios de seleção
% =====================================================================
\section{Área de abrangência e critérios de seleção}
\label{sec:area}

A ação abrange todo o território do Estado do Paraná, contemplando as
Microrregiões de Água e Esgoto. Isso não significa que todos os municípios
serão atendidos. Os territórios e as famílias beneficiadas serão escolhidos
pela aplicação de critérios definidos em âmbito estadual, que variam conforme
a frente de trabalho \cite{mop2026}.

Para as comunidades que receberão sistemas coletivos de abastecimento de água
e, na mesma área, sistemas descentralizados de esgotamento sanitário, o
primeiro critério é o adensamento populacional. Serão consideradas as
comunidades, ou os conjuntos de comunidades, com 70 ou mais domicílios
ocupados, número a partir do qual se espera viabilidade técnica e econômica
compatível com os modelos de gestão a serem estabelecidos. O segundo critério
é o nível de desenvolvimento do município, medido pelo Índice Ipardes de
Desempenho Municipal -- IPDM \cite{ipdm}. O terceiro é a existência de
associações de moradores já constituídas para a gestão de sistemas de
abastecimento de água que necessitem de complementação ou melhorias, o que
favorece a reabilitação de sistemas em funcionamento e aproveita a
organização comunitária existente.

Para as áreas que receberão apenas sistemas descentralizados de esgotamento
sanitário, os critérios estão ligados à proteção dos mananciais. Serão
priorizados os mananciais prioritários do PSH-PR e os mananciais de
abastecimento público selecionados por critérios quali-quantitativos em que o
IDR-Paraná atuará, considerando os domicílios situados em área rural segundo o
Censo Demográfico 2022 \cite{ibge2022}.

A \autoref{fig:selecao} mostra como os dois conjuntos de critérios levam às
duas frentes de implantação descritas nas \autoref{sec:abastecimento} e
\ref{sec:esgotamento}.

\begin{figure}[htbp]
  \centering
  \caption{Aplicação dos critérios de seleção e definição das frentes de trabalho}
  \label{fig:selecao}
  \fluxograma{fluxograma-selecao}
  \fonte{Elaborado pelo IAT (2026) com base em \citeonline{mop2026}.}
\end{figure}

\pendencia{Os critérios não indicam pesos, limiares (exceto os 70
domicílios) nem a ordem de aplicação. Para a reprodutibilidade da seleção,
convém descrever o procedimento de ponderação ou hierarquização e o
responsável por aplicá-lo. O POA 2026--2027 prevê como produto dessa etapa
um mapa dos municípios selecionados (Anexo~\ref{an:cronograma}).}

% =====================================================================
%  3. Arranjo institucional e visão geral da implementação
% =====================================================================
\section{Arranjo institucional e visão geral da implementação}
\label{sec:arranjo}

A ação será executada pelo Instituto Água e Terra (IAT), com coexecução do
Instituto de Desenvolvimento Rural do Paraná (IDR-Paraná) e da Companhia de
Saneamento do Paraná (Sanepar). O IAT responde pelo planejamento, pela
coordenação técnica, pelas contratações, pela fiscalização e pela validação
dos produtos. O IDR-Paraná apoia a mobilização social e o planejamento,
elabora os Planos Integrados de Propriedade (PIP) e atua diretamente na
frente de esgotamento sanitário das áreas de mananciais. A Sanepar executa,
com recursos próprios de contrapartida, o Programa Sanepar Rural. A
Secretaria das Cidades (SECID) participa do planejamento e da coordenação
técnica e valida a implantação dos modelos de gestão. Os municípios integram a
governança e fazem a articulação local, enquanto as comunidades beneficiadas
participam da gestão e da operação cotidiana dos sistemas \cite{mop2026}.

As contratações financiadas pelo Banco Mundial (BIRD) somam custo estimado de
US\$ 35,48 milhões, a serem executados em seis anos. A Sanepar participará
com contrapartida estimada em US\$ 12,31 milhões, destinada ao Programa
Sanepar Rural, que é executado por convênio com os municípios. As demais
atividades serão executadas por contratação pública, mediante licitação,
segundo as regras do Banco Mundial, ou diretamente pela equipe técnica do
IAT, no caso das atividades de planejamento, seleção e articulação
institucional \cite{mop2026}.

Toda contratação financiada pelo BIRD percorre o mesmo ciclo, representado na
\autoref{fig:ciclo-contratacao}. A equipe do IAT elabora o Termo de Referência
(TR) em cerca de três meses, até que ele esteja pronto para envio ao Banco
Mundial. Segue-se o processo licitatório, que se estende até o mês anterior
ao início da execução, quando o contrato é assinado e o serviço, a obra ou o
fornecimento começa. É esse ciclo que estrutura o cronograma de 2026 e 2027
apresentado no Anexo~\ref{an:cronograma}.

\begin{figure}[htbp]
  \centering
  \caption{Ciclo das contratações financiadas pelo BIRD}
  \label{fig:ciclo-contratacao}
  \fluxograma{ciclo-contratacao}
  \fonte{Elaborado pelo IAT (2026) com base no POA 2026--2027.}
\end{figure}

Na execução, as atividades se organizam em três blocos, que a
\autoref{fig:fluxograma-geral} apresenta de forma integrada. O primeiro bloco
é de estruturação institucional e social: a seleção dos municípios e
comunidades, a definição e a pactuação do modelo de gestão e o trabalho
técnico-social junto às famílias. O segundo bloco é a implantação física, que
se divide em duas frentes: o abastecimento de água, que inclui o Programa
Sanepar Rural, e o esgotamento sanitário, que inclui a frente do IDR-Paraná. O
terceiro bloco compreende a entrega dos sistemas, a capacitação dos
beneficiários, o monitoramento e o encerramento. Duas premissas orientam essa
sequência. A primeira é que as comunidades contempladas serão definidas por
meio do trabalho técnico-social. A segunda é que o modelo de gestão será
pactuado com as comunidades antes da instalação dos sistemas, de modo que
nenhuma obra seja entregue sem uma entidade preparada para operá-la
\cite{mop2026}. As seções seguintes descrevem, do começo ao fim, primeiro a
implementação do abastecimento de água e depois a do esgotamento sanitário.

\begin{figure}[htbp]
  \centering
  \caption{Visão geral da implementação do Subcomponente 3.2.b}
  \label{fig:fluxograma-geral}
  \fluxograma{fluxograma-geral}
  \fonte{Elaborado pelo IAT (2026) com base em \citeonline{mop2026}.}
\end{figure}

\pendencia{Os US\$ 35,48 milhões citados no MOP correspondem ao financiamento
BIRD de todo o Subcomponente~3.2, incluindo a política e o plano estadual
(3.2.a, R\$ 10,6 milhões). Pela planilha de custos, o financiamento do 3.2.b
é de R\$ 183,0 milhões, cerca de US\$ 33,5 milhões ao câmbio de R\$ 5,46
\cite{custo2026}. Corrigir o valor.}

\pendencia{O MOP atribui ao IDR-Paraná corresponsabilidade pelas
atividades de abastecimento (Quadro 30: ``IAT/IDR''), mas o texto não descreve
o papel do instituto nessa frente além da mobilização social. Detalhar.}

% =====================================================================
%  4. Implementação do abastecimento de água
%  (Atividades 3.2.2.1, 3.2.2.2, 3.2.2.3 e 3.2.2.6 do MOP)
% =====================================================================
\section{Implementação dos sistemas de abastecimento de água}
\label{sec:abastecimento}

A implementação do abastecimento de água percorre seis fases, desde a
preparação institucional até a operação dos sistemas pela entidade gestora
(\autoref{fig:abastecimento}). Em cada comunidade selecionada, o IAT
pactua com o município e com os moradores o modelo de gestão e realiza o
trabalho técnico-social. Em seguida, perfura e outorga o poço, projeta e
licencia o sistema, executa as obras e, por fim, entrega a infraestrutura a
uma entidade capacitada para operá-la. Essa sequência reúne as
\atividade{3.2.2.1} (modelos de gestão), \atividade{3.2.2.2} (trabalho
técnico-social) e \atividade{3.2.2.3} (sistemas coletivos de abastecimento)
do MOP. Ao final da seção, descreve-se a frente complementar do Programa
Sanepar Rural (\atividade{3.2.2.6}), que segue um arranjo próprio.

\begin{figure}[htbp]
  \centering
  \caption{Implementação do abastecimento de água, do começo ao fim}
  \label{fig:abastecimento}
  \fluxograma{fluxograma-abastecimento}
  \fonte{Elaborado pelo IAT (2026) com base em \citeonline{mop2026} e no POA 2026--2027.}
\end{figure}

% ---------------------------------------------------------------------
\subsection{Seleção e adesão dos municípios e comunidades}
\label{subsec:adesao}

A implementação começa dentro do próprio IAT. Entre maio e outubro de 2026,
a equipe técnica do instituto, em conjunto com o IDR-Paraná, fará o
planejamento integrado das ações de saneamento rural a serem executadas a
partir de 2027. Em paralelo, aplicará os critérios de seleção descritos na
\autoref{sec:area} para identificar os municípios elegíveis, cujo resultado
será consolidado em um mapa. Com esse mapa em mãos, o IAT iniciará a
sensibilização das prefeituras por meio de reuniões e ações institucionais,
apresentando os objetivos do Programa e as responsabilidades de cada ente
\cite{mop2026}.

Os municípios interessados formalizarão sua adesão, o que resultará na lista
de municípios participantes, e assinarão com o IAT um Termo de Cooperação.
Esse termo formaliza a parceria: nele ficam registradas as obrigações do Estado, do município e da comunidade durante
toda a implantação, e ele só será encerrado depois que os sistemas forem
entregues e aceitos pela entidade gestora
(\autoref{subsec:entrega}). Assinado o termo, o IAT fará o contato inicial com
as comunidades, que manifestarão seu interesse por meio de contratos de
pré-adesão. As etapas de seleção, adesão e assinatura do termo estão
previstas no POA para o período de maio a dezembro de 2026, sob
responsabilidade da equipe interna do IAT (Anexo~\ref{an:cronograma}). A
\autoref{fig:adesao} mostra como essas etapas se distribuem entre o IAT, o
município e a comunidade.

\begin{figure}[htbp]
  \centering
  \caption{Seleção e adesão dos municípios e comunidades, por responsável}
  \label{fig:adesao}
  \fluxograma{fluxograma-adesao}
  \fonte{Elaborado pelo IAT (2026) com base em \citeonline{mop2026} e no POA 2026--2027.}
\end{figure}

% ---------------------------------------------------------------------
\subsection{Definição e implantação do modelo de gestão}
\label{subsec:modelos-gestao}

O modelo de gestão é o que garante que o sistema continuará funcionando
depois da entrega. A dispersão das comunidades rurais e a limitada capacidade
operacional de muitos municípios tornam necessários arranjos institucionais
que promovam a cooperação entre os diferentes níveis de governo. Os modelos
específicos para o saneamento rural ainda estão em definição. A alternativa
preliminar em análise é um modelo baseado em arranjos intermunicipais,
possivelmente estruturado por meio de consórcios públicos ou mecanismos
equivalentes, que fortaleça a capacidade de gestão dos municípios e garanta
suporte técnico às comunidades. Nesse modelo, Estado, municípios e comunidades
compartilham responsabilidades: o IAT, a SECID e o IDR-Paraná apoiam o
planejamento, a coordenação técnica, a mobilização social e o monitoramento;
os municípios participam da governança e da articulação local; e as
comunidades participam da gestão e da operação cotidiana dos sistemas
\cite{mop2026}. Segundo o POA, a articulação com os entes envolvidos para
definir um modelo de gestão piloto ocorrerá entre abril e dezembro de 2026,
conduzida pelo IAT com o IDR-Paraná e a Sanepar.

Uma vez definido o modelo, sua implantação nas comunidades seguirá seis
etapas, representadas na \autoref{fig:modelo-gestao}. Nos estudos
preliminares, entre 2027 e 2028, serão selecionados os instrumentos de
gestão, definidos os parceiros, critérios e indicadores, escolhidos os meios
de comunicação do modelo e feitos os primeiros contatos com o município e a
comunidade, até a assinatura dos contratos de pré-adesão. Na etapa de
implantação, entre 2028 e 2029, serão elaborados o estatuto, os manuais e as
cartilhas, feito o enquadramento jurídico da proposta, colhida a adesão de
cada família e regularizada a associação que assumirá a gestão. A etapa de
gestão propriamente dita, entre 2030 e 2032, compreende a elaboração dos
manuais de desempenho operacional, os investimentos necessários, a
estruturação do monitoramento e o treinamento das pessoas que vão operar o
sistema. Ao longo de todo esse período, a implantação será acompanhada por
supervisão técnica e avaliada semestralmente, com os dados registrados no
SIGARH e no GeoPR. A atividade se encerra entre 2032 e 2033, com a
consolidação dos resultados, os relatórios finais e a prestação de contas.
Para dar suporte a essa estrutura, o POA prevê ainda a aquisição de
infraestrutura para implantação do modelo de gestão, cujo TR começará a ser
elaborado em novembro de 2027.

\begin{figure}[htbp]
  \centering
  \caption{Etapas de implantação do modelo de gestão e seus produtos}
  \label{fig:modelo-gestao}
  \fluxograma{fluxograma-modelo-gestao}
  \fonte{Elaborado pelo IAT (2026) com base em \citeonline[Quadro 28]{mop2026}.}
\end{figure}

O sucesso desta atividade será medido pelo número de associações de
saneamento rural com modelo de gestão adotado, indicador de resultado do
Componente~3 cuja meta é de 122 associações até 2032
(\autoref{tab:metas-componente3}). O detalhamento das etapas, com os
responsáveis, os produtos, as evidências e a validação de cada uma, está no
\autoref{qua:etapas-3221}, no Apêndice~\ref{ap:quadros}.

\pendencia{No Quadro 28 do MOP (\autoref{qua:etapas-3221}), a etapa
``Supervisão técnica'' repete o produto e as evidências das obras
(``supervisão e gerenciamento de obras'', ``laudos de vistoria e testes de
estanqueidade''). A sigla ``EC'' não está definida. Além disso, o Quadro 31
adota ``tarifa instituída'' como evidência e a SECID como validadora, enquanto
o Quadro 28 usa associações regularizadas e pessoas capacitadas, validadas
pelo IAT. Harmonizar.}

% ---------------------------------------------------------------------
\subsection{Mobilização social e trabalho técnico-social}
\label{subsec:tts}

Com o município e a comunidade formalmente engajados, inicia-se o trabalho
técnico-social, que tem três finalidades: conhecer as famílias e suas
condições de saneamento, definir quem será atendido e validar socialmente as
soluções técnicas propostas. Esse trabalho será executado por empresa
contratada segundo as regras do BIRD. No POA, a mesma contratação reúne a
mobilização social, o levantamento preliminar social, o cadastro das famílias
por busca ativa e os estudos técnicos iniciais do poço (estudo
hidrogeológico, anuência prévia, teste de vazão, análise físico-química e
anteprojeto), descritos na \autoref{subsec:captacao}. O TR dessa
contratação será elaborado entre setembro e novembro de 2026 e a licitação
se estenderá até setembro de 2027, de modo que a execução em campo começará
em outubro de 2027.

Em campo, a empresa começará pela mobilização e sensibilização dos
municípios, das comunidades, das associações e dos proprietários, informando
a população sobre o Programa, estimulando a participação comunitária e
reconhecendo o território onde os sistemas serão implantados. Em seguida,
fará a busca ativa das famílias, inclusive das que vivem em áreas remotas,
acompanhada de ações de capacitação e organização comunitária. Com base
nesse contato, será feito um diagnóstico coletivo das formas atuais de
abastecimento de água e de esgotamento sanitário, que avaliará as condições
existentes, a demanda de água, a infraestrutura necessária e a viabilidade
dos sistemas, complementado pelo levantamento das condições de cada
domicílio.

O diagnóstico permite aplicar os critérios de enquadramento, que combinam
aspectos sociais, territoriais e técnicos, para definir quais famílias serão
atendidas. Cada família enquadrada assinará um termo de adesão, no qual
declara sua participação e seu compromisso com o uso e a manutenção
adequados do sistema. Os resultados serão compilados e apresentados em um
relatório final, que constitui a evidência de conclusão da atividade e dá
baixa na etapa correspondente do Termo de Cooperação. Esse diagnóstico tem
duplo uso: orienta a captação e o projeto do sistema de água e identifica,
domicílio a domicílio, a situação do esgoto, informação de que depende a
frente de esgotamento (\autoref{sec:esgotamento}). A
\autoref{fig:tts} resume a sequência e os produtos.

\begin{figure}[htbp]
  \centering
  \caption{Sequência e produtos do trabalho técnico-social}
  \label{fig:tts}
  \fluxograma{fluxograma-tts}
  \fonte{Elaborado pelo IAT (2026) com base em \citeonline{mop2026} e no POA 2026--2027.}
\end{figure}

\pendencia{O MOP separa o trabalho técnico-social (\atividade{3.2.2.2}) dos
estudos do poço (\atividade{3.2.2.3}), mas o POA os reúne em uma única
contratação. O MOP também coloca a contratação da empresa depois do
diagnóstico, em um item cuja descrição abrange projetos e obras. Harmonizar o
escopo e a sequência.}

% ---------------------------------------------------------------------
\subsection{Captação, projeto e licenciamento}
\label{subsec:captacao}

Definidas as comunidades e as famílias, começam os estudos de engenharia,
que seguem duas linhas encadeadas (\autoref{fig:captacao-projeto}). A primeira
trata da captação subterrânea. Os estudos preliminares avaliarão a
viabilidade técnica, social e ambiental do sistema e verificarão a
titularidade e a situação fundiária das áreas onde ficarão o poço e o
reservatório, cuidado que evita conflitos posteriores sobre o uso da terra.
O levantamento de campo reunirá os dados geológicos e hidrogeológicos
necessários para identificar os aquíferos e as condições ambientais do
local. Com base nele, será definida a locação do ponto de captação e
elaborado o projeto construtivo do poço, que dependerá de anuência prévia
antes da perfuração. Perfurado o poço, serão realizados o teste de vazão e a
análise físico-química da água, que confirmarão se há quantidade e qualidade
suficientes para o abastecimento. Por fim, serão feitos os estudos ambientais
e sociais exigidos para a obtenção da outorga de direito de uso da água
\cite{mop2026}. Para essa linha, o POA prevê a aquisição dos materiais de
perfuração, com TR entre outubro e dezembro de 2026, licitação ao longo de
2027 e início da execução em novembro de 2027.

A segunda linha trata do sistema de tratamento, reservação e distribuição. A
vazão e a qualidade da água obtidas no poço orientam o anteprojeto da
solução de tratamento e armazenamento. Após a ordenação das despesas e a
contratação, serão elaborados o projeto básico e o projeto executivo da casa
química, do reservatório, da parte eletromecânica e da rede de distribuição,
que precisarão ser aprovados antes da obra. Paralelamente, serão realizados
os estudos ambientais e sociais e obtidas as licenças ambientais. O POA prevê
a contratação dos projetos das redes de água e de esgoto com TR entre
novembro de 2026 e janeiro de 2027 e início da execução em dezembro de 2027.
Somente com a outorga e as licenças obtidas as obras poderão começar. O
\autoref{qua:etapas-3223}, no Apêndice~\ref{ap:quadros}, detalha essas etapas
com seus produtos e evidências.

\begin{figure}[htbp]
  \centering
  \caption{Captação subterrânea, projeto e licenciamento do sistema de abastecimento}
  \label{fig:captacao-projeto}
  \fluxograma{fluxograma-captacao-projeto}
  \fonte{Elaborado pelo IAT (2026) com base em \citeonline{mop2026}.}
\end{figure}

\pendencia{O MOP prevê ``contratação integrada'' para o poço e para o sistema
e a define como aquela em que o contratado elabora o projeto executivo e
executa a obra. Na Lei nº~14.133/2021 (art.~6º, XXXII e XXXIII), essa
definição corresponde à contratação \textit{semi-integrada}. Já o POA prevê
outro arranjo: aquisição de materiais pelo IAT e contratação separada dos
projetos. Definir o modelo de execução, indicar quem executará as obras com
os materiais adquiridos e verificar a compatibilidade com o Regulamento de
Aquisições do Banco Mundial.}

% ---------------------------------------------------------------------
\subsection{Execução das obras, supervisão e fiscalização}
\label{subsec:obras}

A obra compreende a construção da casa química, do reservatório, das
instalações eletromecânicas e da rede de distribuição, com as ligações
domiciliares e os hidrômetros. Para essa fase, o POA prevê a aquisição de
materiais para a implantação das redes de abastecimento em dois lotes, de 25
e de 30 unidades. O primeiro terá TR entre dezembro de 2026 e fevereiro de
2027 e licitação até dezembro de 2027; o TR do segundo começará em dezembro
de 2027.

Antes do início das obras, o IAT contratará uma empresa de supervisão e
gerenciamento, que fará o acompanhamento técnico independente da execução.
O mesmo contrato atenderá também às obras de esgotamento sanitário, e seu TR
será elaborado entre dezembro de 2026 e fevereiro de 2027, segundo o POA. A
execução é acompanhada em três níveis, como mostra a \autoref{fig:obras}. O
executor constrói a infraestrutura conforme o projeto. A supervisão
acompanha a obra e registra, em relatórios periódicos, o andamento, os
problemas encontrados e as recomendações. A fiscalização do IAT verifica se
a obra segue o projeto e as normas técnicas e registra essa verificação em
relatórios de fiscalização. Quando a verificação apontar não conformidades,
o executor fará os ajustes necessários antes de nova verificação. Estando
tudo conforme, o IAT emitirá o atesto de recebimento, acompanhado dos laudos
de vistoria, e o executor entregará o projeto \textit{as built}
\cite{mop2026}.

\begin{figure}[htbp]
  \centering
  \caption{Execução, supervisão e fiscalização das obras, por responsável}
  \label{fig:obras}
  \fluxograma{fluxograma-obras}
  \fonte{Elaborado pelo IAT (2026) com base em \citeonline{mop2026}.}
\end{figure}

\pendencia{Os lotes de materiais previstos no POA (25 + 30 = 55 unidades)
não cobrem os 83 sistemas da meta. Confirmar se há lotes posteriores a
2027. A planilha de custos prevê a implantação dos 83 sistemas entre 2028 e
2033; os prazos dos Quadros 29 e 30 do MOP devem ser ajustados a ela. A
quantidade de hidrômetros ainda está como ``XXX''.}

% ---------------------------------------------------------------------
\subsection{Pré-operação, capacitação e entrega à entidade gestora}
\label{subsec:entrega}

Concluída a obra, o sistema passará por um período de pré-operação e
testes, durante o qual os operadores locais serão capacitados. Os
beneficiários também receberão treinamentos sobre operação e manutenção,
comprovados por listas de presença e certificados. A transferência será
formalizada em uma reunião técnica com a comunidade e a entidade gestora, na
qual o IAT entregará o sistema e a entidade confirmará formalmente o seu
recebimento. A partir desse momento, a entidade gestora assume a operação e a
manutenção, enquanto o IAT mantém o suporte técnico pós-implantação e avalia
semestralmente o desempenho dos sistemas por meio de relatórios de
conformidade técnica e ambiental registrados no SIGARH e no GeoPR. A
atividade se encerra com a consolidação dos resultados, o relatório final, a
emissão do Termo de Recebimento Definitivo (TRD) e o encerramento do Termo de
Cooperação com o município \cite{mop2026}. A \autoref{fig:entrega} representa
essa transição.

\begin{figure}[htbp]
  \centering
  \caption{Transferência, operação e encerramento dos sistemas de abastecimento}
  \label{fig:entrega}
  \fluxograma{fluxograma-entrega}
  \fonte{Elaborado pelo IAT (2026) com base em \citeonline{mop2026}.}
\end{figure}

% ---------------------------------------------------------------------
\subsection{Programa Sanepar Rural}
\label{subsec:sanepar-rural}

Nos municípios que mantêm Contrato de Concessão ou de Programa com a Sanepar,
o abastecimento das comunidades rurais será ampliado pelo Programa Sanepar
Rural, com recursos de contrapartida da companhia e execução por convênio
\cite{mop2026}. Estão previstos 222 sistemas coletivos entre 2027 e 2031, com
investimento de R\$ 67,1 milhões \cite{custo2026}. Nesse arranjo, a Sanepar fornece a expertise, os estudos técnicos, o apoio técnico, ambiental e
sociocomunitário, o treinamento, os materiais hidráulicos e os equipamentos,
e cabe ao município executar as obras, com o apoio técnico da companhia. As atividades
seguem o manual do programa \cite{saneparrural}. Começam com o levantamento
técnico e o estudo de engenharia e seguem com a formalização da parceria, por
meio de Termo de Compromisso e Responsabilidade vinculado ao contrato de
concessão. Depois vêm a elaboração dos projetos de engenharia do sistema de
abastecimento, o fornecimento dos materiais e a execução das obras pelo
município. Paralelamente, a comunidade é mobilizada e apoiada na formação de
sua associação de moradores e capacitada em uso racional da água, higiene e
manutenção preventiva. Após a obra, a Sanepar acompanha a qualidade da água e
presta suporte operacional às associações (\autoref{fig:sanepar-rural}). Os
modelos de gestão descritos na \autoref{subsec:modelos-gestao} também se
aplicam a esses sistemas.

\begin{figure}[htbp]
  \centering
  \caption{Programa Sanepar Rural, por responsável}
  \label{fig:sanepar-rural}
  \fluxograma{fluxograma-sanepar-rural}
  \fonte{Elaborado pelo IAT (2026) com base em \citeonline{mop2026}.}
\end{figure}

\pendencia{A meta do Programa Sanepar Rural no Quadro 31 do MOP (59.780
pessoas) é a meta total do indicador ``Pessoas atendidas com água gerida de
forma segura'' do Componente~3, que soma os sistemas do IAT (cerca de 16.200
pessoas) e da Sanepar (cerca de 43.500). A parcela da Sanepar Rural é de
cerca de 43.500 pessoas. O MOP também indica periodicidade ``anual'' no texto e
``contínua'' no Quadro 31, e o manual do programa deve ser citado com título,
versão e data.}

O \autoref{qua:sintese-abastecimento} resume, para cada atividade do
abastecimento, o responsável, o prazo, o produto e a forma de comprovação.

\begin{quadro}[htbp]
\caption{Síntese das atividades de abastecimento de água}
\label{qua:sintese-abastecimento}
\footnotesize
\renewcommand{\arraystretch}{1.25}
\begin{tabularx}{\textwidth}{|P{2.9cm}|P{2.0cm}|Q{1.3cm}|L|L|P{1.4cm}|}
\hline
\cab{Atividade} & \cab{Responsável} & \cab{Prazo} & \cab{Produto} &
\cab{Evidência} & \cab{Validação} \\ \hline
3.2.2.1 Modelos de gestão & IAT, IDR-Paraná e Sanepar & 2027 a 2033 &
Gestão dos sistemas implantada & Tarifa instituída & SECID \\ \hline
3.2.2.2 Trabalho técnico-social & IAT e IDR-Paraná & 2027 a 2030 &
Diagnóstico e cadastro domiciliar & Relatório final e baixa no Termo de
Cooperação & IAT e IDR-Paraná \\ \hline
3.2.2.3 Sistemas coletivos de abastecimento & IAT e IDR-Paraná & 2028 a 2033 &
83 sistemas implantados & Termo de Recebimento Definitivo & IAT \\ \hline
3.2.2.6 Programa Sanepar Rural & Sanepar e IAT & 2027 a 2031 &
222 sistemas implantados (pessoas atendidas) & Relatório técnico aprovado &
Sanepar \\ \hline
\end{tabularx}
\fonte{Adaptado de \citeonline[Quadros 30 e 31]{mop2026}; prazos conforme \citeonline{custo2026}.}
\end{quadro}

% =====================================================================
%  5. Implementação do esgotamento sanitário
%  (Atividades 3.2.2.4 e 3.2.2.5 do MOP)
% =====================================================================
\section{Implementação dos sistemas de esgotamento sanitário}
\label{sec:esgotamento}

O esgotamento sanitário será resolvido por sistemas descentralizados, isto é,
soluções individuais instaladas em cada domicílio, adequadas à baixa
densidade das áreas rurais. A implementação tem duas portas de entrada
(\autoref{fig:esgotamento-geral}). A frente das comunidades atende 70\% dos
domicílios das comunidades que recebem o sistema coletivo de abastecimento
de água, o que corresponde a 4.067 domicílios (\atividade{3.2.2.4}). A
frente dos mananciais atende 2.300 domicílios em propriedades rurais
situadas nas áreas prioritárias do PSH-PR em que atua o IDR-Paraná
(\atividade{3.2.2.5}). As duas frentes diferem na forma como os domicílios
são identificados e contratados, mas, a partir do projeto de cada
propriedade, seguem o mesmo procedimento de instalação, verificação,
capacitação, medição e pagamento \cite{mop2026}.

\begin{figure}[htbp]
  \centering
  \caption{Implementação do esgotamento sanitário descentralizado, do começo ao fim}
  \label{fig:esgotamento-geral}
  \fluxograma{fluxograma-esgotamento-geral}
  \fonte{Elaborado pelo IAT (2026) com base em \citeonline{mop2026} e no POA 2026--2027.}
\end{figure}

% ---------------------------------------------------------------------
\subsection{Identificação dos domicílios e definição das soluções}
\label{subsec:esgotamento-demanda}

Na frente das comunidades, o ponto de partida é o trabalho técnico-social
descrito na \autoref{subsec:tts}. Depois da seleção e da mobilização das
prefeituras e das comunidades, o diagnóstico domiciliar identificará como
cada família dispõe hoje o seu esgoto. É a partir dessa informação que se
definem os domicílios a serem atendidos \cite{mop2026}. Antes de contratar a
instalação, o IAT, com o IDR-Paraná, fará um estudo para definir as melhores
alternativas de sistemas de esgotamento sanitário para comunidades rurais.
Esse estudo será conduzido pela equipe interna entre outubro e dezembro de
2026, segundo o POA, e servirá de base técnica para os termos de referência.

Na frente dos mananciais, a demanda nasce dos Planos Integrados de Propriedade
(PIP) elaborados pelo IDR-Paraná no âmbito do Subcomponente 2.3. Concluídos os
PIP, será iniciado o processo de instalação dos sistemas nas propriedades
situadas nos mananciais prioritários do PSH-PR \cite{mop2026}.

\pendencia{O MOP não indica a tecnologia de tratamento prevista nem as normas
técnicas de projeto. Depois do estudo de alternativas, citar as normas
aplicáveis, como a ABNT NBR 7229 (tanques sépticos) e a ABNT NBR 13969
(unidades de tratamento complementar e disposição final).}

% ---------------------------------------------------------------------
\subsection{Contratação}
\label{subsec:esgotamento-contratacao}

Em ambas as frentes, o IAT elaborará os termos de referência para a
aquisição das soluções descentralizadas e para a contratação dos serviços de
obras civis de instalação. Aprovados os TR, serão abertos os editais,
realizadas as licitações e contratadas as empresas \cite{mop2026}. O POA
distingue duas contratações. Na frente das comunidades, a aquisição de
material para implantação do esgotamento terá TR entre janeiro e março de
2027 e licitação de abril a dezembro de 2027. Na frente do IDR-Paraná, a
aquisição de materiais e a implantação dos sistemas serão contratadas em
conjunto, com TR entre abril e junho de 2027 e licitação de julho a dezembro de
2027. A supervisão dessas obras ficará a cargo da mesma empresa contratada
para supervisionar as obras de abastecimento (\autoref{subsec:obras}).

\pendencia{O POA fala em ``rede de esgotamento sanitário'' na contratação da
frente das comunidades, enquanto o MOP prevê sistemas descentralizados
individuais. Uniformizar a terminologia.}

% ---------------------------------------------------------------------
\subsection{Projeto, instalação e verificação técnica}
\label{subsec:esgotamento-instalacao}

Contratada a empresa, ela elaborará um projeto técnico específico para cada
residência, considerando as características do terreno, o número de moradores
e a solução de tratamento mais adequada. Em seguida, executará as obras civis
necessárias à instalação das estruturas e do dispositivo de esgotamento
domiciliar. Concomitantemente à instalação, os proprietários serão orientados
sobre o funcionamento do sistema, as práticas adequadas de uso e a manutenção
preventiva, com registro dessa capacitação \cite{mop2026}.

Concluída a instalação, será feita uma verificação técnica para atestar que o
sistema funciona conforme os requisitos definidos e para registrar a sua
garantia. Essa verificação funciona como um ponto de decisão. Se o sistema
atender às condições técnicas, é validado e segue para a medição. Se estiver
em desacordo, a empresa fará as correções necessárias e o sistema será
verificado novamente, até que seja aprovado (\autoref{fig:esgotamento-instalacao}).
Os testes de estanqueidade após a instalação, previstos entre as medidas
socioambientais do Subcomponente~3.2, reduzem o risco de contaminação do solo
e do lençol freático por instalação incorreta.

\begin{figure}[htbp]
  \centering
  \caption{Projeto, instalação e verificação técnica dos sistemas descentralizados}
  \label{fig:esgotamento-instalacao}
  \fluxograma{fluxograma-esgotamento}
  \fonte{Elaborado pelo IAT (2026) com base em \citeonline{mop2026}.}
\end{figure}

% ---------------------------------------------------------------------
\subsection{Medição, comprovação e pagamento}
\label{subsec:esgotamento-medicao}

Como os sistemas são instalados propriedade por propriedade, o pagamento
depende de uma cadeia de comprovação que envolve a empresa, o proprietário e
o IAT (\autoref{fig:esgotamento-medicao}). A empresa registrará, em relatórios
de medição, o andamento das atividades e as intervenções executadas em cada
propriedade, e reunirá a documentação que comprova a elaboração e a execução
dos projetos, com as evidências das intervenções. O proprietário, por sua
vez, assinará o registro que confirma a execução das melhorias em sua
propriedade \cite{mop2026}.

Com esse material, o IAT fará uma auditoria documental e visitas técnicas de
verificação \textit{in loco} para confirmar a execução das intervenções. A
empresa consolidará então as informações no relatório de medição final, e o
IAT avaliará por amostragem, também \textit{in loco}, a qualidade das unidades
atendidas. Só depois dessas verificações serão realizadas a ordenação das
despesas e o pagamento dos serviços, seguidos do encerramento formal do
processo. O relatório técnico aprovado e o comprovante de pagamento
constituem a evidência de conclusão de cada lote \cite{mop2026}.

\begin{figure}[htbp]
  \centering
  \caption{Medição, comprovação, verificação independente e pagamento, por responsável}
  \label{fig:esgotamento-medicao}
  \fluxograma{fluxograma-esgotamento-medicao}
  \fonte{Elaborado pelo IAT (2026) com base em \citeonline{mop2026}.}
\end{figure}

\pendencia{O MOP não define o tamanho da amostra da verificação \textit{in
loco} nem a diferença entre ``verificação \textit{in loco}'' e ``amostragem
\textit{in loco}''. Também não indica quem executa cada verificação; nesta
estrutura, elas foram atribuídas ao IAT, responsável pela validação.
Confirmar.}

% ---------------------------------------------------------------------
\subsection{Particularidades da frente do IDR-Paraná}
\label{subsec:esgotamento-idr}

A frente do IDR-Paraná segue o mesmo procedimento de projeto, instalação,
verificação, capacitação, medição e pagamento descrito acima. As diferenças
estão na origem da demanda, no território e nos responsáveis
(\autoref{qua:diferencas-esgotamento}). Os domicílios são identificados pelos
PIP e não pelo trabalho técnico-social. O território é o dos mananciais de
abastecimento público, onde o esgoto doméstico não tratado representa risco
direto para a água que abastece as cidades. O projeto é chamado de projeto
sanitário doméstico da propriedade. Por fim, a validação é compartilhada entre
o IAT e o IDR-Paraná e não há dependência das prefeituras \cite{mop2026}.

\begin{quadro}[htbp]
\caption{Diferenças entre as frentes de esgotamento sanitário}
\label{qua:diferencas-esgotamento}
\small
\renewcommand{\arraystretch}{1.25}
\begin{tabularx}{\textwidth}{|P{3.2cm}|L|L|}
\hline
\cab{Aspecto} & \cab{Frente comunidades (3.2.2.4)} & \cab{Frente mananciais -- IDR-Paraná (3.2.2.5)} \\ \hline
Origem da demanda & Diagnóstico do trabalho técnico-social &
  Planos Integrados de Propriedade (PIP) do IDR-Paraná \\ \hline
Território & Comunidades atendidas com abastecimento de água &
  Mananciais prioritários do PSH-PR \\ \hline
Meta & 4.067 domicílios (70\% dos atendidos com água) & 2.300 domicílios \\ \hline
Contratação (POA) & Aquisição de material; TR jan--mar/2027 &
  Aquisição de materiais e implantação; TR abr--jun/2027 \\ \hline
Prazo (planilha de custos) & 2028 a 2033 & 2029 a 2033 \\ \hline
Dependências & Prefeituras & -- \\ \hline
Evidência & \multicolumn{2}{l|}{Relatório técnico aprovado e comprovante de pagamento, por lote} \\ \hline
Validação & IAT & IAT e IDR-Paraná \\ \hline
\end{tabularx}
\fonte{Adaptado de \citeonline[Quadros 30 e 31]{mop2026} e do POA 2026--2027.}
\end{quadro}

% =====================================================================
%  6. Salvaguardas ambientais e sociais
% =====================================================================
\section{Salvaguardas ambientais e sociais}
\label{sec:salvaguardas}

Todas as fases descritas nas seções anteriores observarão o Quadro Ambiental
e Social do Banco Mundial \cite{bm-esf}, aplicado ao PSH-PR por meio de cinco
instrumentos. O Marco de Gestão Ambiental e Social (MGAS) orienta o controle
dos impactos das obras e da operação. O Plano de Engajamento das Partes
Interessadas (PEPI) orienta a participação das comunidades. Os Procedimentos
de Gestão de Mão de Obra (PGR) tratam da saúde e da segurança dos
trabalhadores. O Marco de Política de Reassentamento (MRF) rege as questões
de posse e uso da terra, e o Mecanismo de Queixas e Reclamações (MQR) recebe e
responde às manifestações da população.

Os riscos variam conforme a fase. No planejamento, o principal risco é
excluir grupos vulneráveis ou comunidades tradicionais dos diagnósticos, o
que será mitigado com consultas públicas e metodologias participativas na
seleção das áreas, além da busca ativa de famílias em áreas remotas e de
reuniões comunitárias em linguagem simples. Na captação e na distribuição de
água, os riscos são a contaminação de aquíferos e os impactos temporários das
obras, como ruído, poeira e erosão, controlados pela adoção de normas
técnicas de perfuração, pelo gerenciamento de resíduos e pela umectação das
vias. Como o poço e o reservatório ocupam áreas específicas, a titularidade
dessas áreas será verificada previamente e formalizada por anuências e termos
de servidão, conforme o MRF. Na operação, os riscos são a falta de
sustentabilidade técnica e os conflitos pelo uso da água, que o modelo de
gestão, a tarifa social e o treinamento das associações procuram evitar. No
esgotamento, os riscos são o manejo inadequado do lodo e a contaminação do
solo por instalação incorreta, enfrentados com testes de estanqueidade e
capacitação dos proprietários. Nas obras em geral, protocolos de saúde e
segurança do trabalho protegerão operários, técnicos e famílias. O MQR será
implantado e divulgado nas comunidades, e seus tempos de resposta serão
monitorados semestralmente \cite[Quadros 25 e 26]{mop2026}. O
\autoref{qua:salvaguardas} organiza esses riscos por atividade.

\begin{quadro}[htbp]
\caption{Riscos ambientais e sociais e medidas de mitigação do Subcomponente 3.2.b}
\label{qua:salvaguardas}
\footnotesize
\renewcommand{\arraystretch}{1.25}
\begin{tabularx}{\textwidth}{|P{2.4cm}|L|L|P{2.5cm}|P{1.9cm}|}
\hline
\cab{Atividade} & \cab{Riscos potenciais} & \cab{Medidas de mitigação} &
\cab{Instrumento} & \cab{Acompanha\-mento} \\ \hline
Planejamento e gestão &
Exclusão de grupos vulneráveis ou comunidades tradicionais nos diagnósticos rurais &
Consultas públicas e metodologias participativas para seleção das áreas &
PEPI & IAT \\ \hline
Sistemas de abastecimento e distribuição &
Contaminação de aquíferos; impactos temporários de obra (ruído, poeira e erosão local) &
Adoção de normas técnicas de perfuração &
MGAS & IAT \\ \hline
Operação e manutenção &
Falta de sustentabilidade técnica; conflitos locais pelo uso da água &
Treinamento das associações locais; definição de modelos de gestão e tarifa social &
MGAS e Plano de Gestão Social & IAT e prefeituras \\ \hline
Sistemas de tratamento de esgoto &
Manejo inadequado de lodo; contaminação do solo por instalação incorreta &
Capacitação das Unidades de Produção Familiar (UPF); observância da RDC 49/2013 &
MGAS e protocolos de assistência técnica & IAT / IDR-Paraná \\ \hline
Relações de trabalho nas obras &
Acidentes de trabalho ou condições inadequadas para operários e técnicos &
Protocolos de saúde e segurança do trabalho &
PGR & IAT \\ \hline
\end{tabularx}
\fonte{Adaptado de \citeonline[Quadro 26]{mop2026}.}
\end{quadro}

\pendencia{Conferir a citação da RDC 49/2013 (Anvisa). Essa resolução trata
da regularização de atividades de interesse sanitário do microempreendedor
individual e do empreendimento familiar rural, e não do manejo de lodo de
sistemas de esgotamento. Verificar qual norma se pretende citar.}

% =====================================================================
%  7. Monitoramento, avaliação e indicadores
% =====================================================================
\section{Monitoramento, avaliação e indicadores}
\label{sec:monitoramento}

O monitoramento acompanha a ação em dois níveis. No nível da execução, os
indicadores do \autoref{qua:indicadores} medem as entregas físicas: poços e
sistemas implantados, domicílios atendidos com água e com esgotamento e
população beneficiada. No nível dos resultados, a ação contribui para três
indicadores do Componente~3 (\autoref{tab:metas-componente3}): pessoas
atendidas com água gerida de forma segura, pessoas atendidas com pelo menos
esgotamento sanitário e associações de saneamento rural com modelo de gestão
adotado. Os valores de pessoas atendidas consideram 2,8 moradores por
domicílio \cite{ibge2022} e, no caso da água, incluem os sistemas implantados
pela Sanepar \cite{mop2026}.

\begin{quadro}[htbp]
\caption{Indicadores de implantação e reabilitação dos sistemas de abastecimento de água e esgotamento sanitário no meio rural}
\label{qua:indicadores}
\footnotesize
\renewcommand{\arraystretch}{1.25}
\begin{tabularx}{\textwidth}{|L|L|Q{1.6cm}|P{2.2cm}|Q{1.4cm}|}
\hline
\cab{Indicador} & \cab{Fórmula de cálculo} & \cab{Horizonte} &
\cab{Unidade} & \cab{Meta final} \\ \hline
Projetos de poços tubulares com sistemas de abastecimento de água implantados &
Quantidade de projetos executados & 1º ao 6º ano & Projetos (nº) & 83 \\ \hline
Domicílios atendidos com sistemas de abastecimento de água &
Soma das comunidades dentro dos municípios contemplados pelas contratações &
1º ao 6º ano & Domicílios & 5.810 \\ \hline
Domicílios atendidos com esgotamento sanitário &
70\% dos domicílios atendidos com abastecimento de água &
4º ao 6º ano & Domicílios & 4.067 \\ \hline
Domicílios atendidos com esgotamento sanitário nas propriedades rurais do IDR-Paraná &
Obtidos a partir dos dados do Censo 2022 & 4º ao 6º ano & Domicílios & 2.300 \\ \hline
População rural atendida com acesso à água potável e ao esgotamento sanitário &
Soma dos habitantes beneficiados pelas obras & 4º ao 6º ano &
Pessoas (nº) & 22.708 \\ \hline
\end{tabularx}
\fonte{Adaptado de \citeonline[Quadro 27]{mop2026}.}
\end{quadro}

\begin{table}[htbp]
\centering
\caption{Metas progressivas dos indicadores do Componente 3 relacionados ao Subcomponente 3.2.b}
\label{tab:metas-componente3}
\footnotesize
\renewcommand{\arraystretch}{1.2}
\begin{adjustbox}{max width=\linewidth}
\begin{tabular}{lrrrrrrr}
\toprule
\textbf{Indicador} & \textbf{2025} & \textbf{2027} & \textbf{2028} &
\textbf{2029} & \textbf{2030} & \textbf{2031} & \textbf{2032} \\
\midrule
Pessoas com água gerida de forma segura & 0 & 980 & 10.900 & 23.000 & 35.400 & 49.000 & 59.780 \\
\quad sexo feminino & 0 & 400 & 5.200 & 11.000 & 16.000 & 23.300 & 28.500 \\
\quad jovens (15 a 24 anos) & 0 & 110 & 1.300 & 2.800 & 4.200 & 5.900 & 7.000 \\
Pessoas com, pelo menos, esgotamento sanitário & 0 & 126 & 3.000 & 6.700 & 10.600 & 14.600 & 17.800 \\
\quad sexo feminino & 0 & 60 & 1.400 & 3.200 & 5.000 & 6.900 & 8.000 \\
\quad jovens (15 a 24 anos) & 0 & 15 & 360 & 800 & 1.200 & 1.700 & 2.000 \\
Associações RWSS com modelo de gestão adotado & 0 & 0 & 0 & 10 & 30 & 90 & 122 \\
\bottomrule
\end{tabular}
\end{adjustbox}
\fonte{Adaptado de \citeonline[Quadro 23]{mop2026}.}
\end{table}

\pendencia{Revisar as fórmulas do \autoref{qua:indicadores}. O indicador de
domicílios com abastecimento é calculado como ``soma das comunidades'', uma
unidade diferente de domicílios. O MOP escreve ``70\% do total de pessoas
atendidas domicílios atendidos com abastecimento'', o que mistura pessoas e
domicílios. O horizonte do esgotamento (4º ao 6º ano) não coincide com o
prazo das atividades (anos 2 a 5). O MOP também remete a um ``Quadro XX''.}

As informações que alimentam esses indicadores percorrem um fluxo definido
(\autoref{fig:monitoramento}). A supervisão e a fiscalização das obras
produzem, de forma contínua, relatórios de supervisão, laudos de vistoria e
testes de estanqueidade, validados pelo IAT. A cada semestre, o IAT
consolida essas informações, junto com o acompanhamento dos tempos de
resposta do Mecanismo de Queixas e Reclamações, em relatórios de
conformidade técnica e ambiental, cujos dados são registrados no SIGARH e no
GeoPR. Os indicadores de resultado têm frequência anual. O IAT e a Sanepar
reportam à Unidade de Gerenciamento do Programa (UGP) o número acumulado de
domicílios conectados e de associações que adotaram o modelo de gestão, a
partir de seus próprios registros e do banco de dados das associações. A UGP
consolida e valida essas informações à luz do MGAS e as reporta ao Banco
Mundial, que também valida, ao final, o relatório de encerramento e o Termo de
Recebimento Definitivo \cite{mop2026}.

\begin{figure}[htbp]
  \centering
  \caption{Fluxo de monitoramento, validação e reporte}
  \label{fig:monitoramento}
  \fluxograma{fluxograma-monitoramento}
  \fonte{Elaborado pelo IAT (2026) com base em \citeonline{mop2026}.}
\end{figure}

% =====================================================================
%  8. Cronograma de execução
% =====================================================================
\section{Cronograma de execução}
\label{sec:cronograma}

O programa tem sete anos de execução, de 2027 a 2033, precedidos de um ano
de preparação institucional em 2026 (\autoref{fig:cronograma}). Os prazos
seguem a planilha de custos do PSH-PR \cite{custo2026}, que é a referência
para as metas anuais (\autoref{tab:metas-anuais}). Os ``Anos 0 a 6'' usados
nos quadros de etapas do MOP (Apêndice~\ref{ap:quadros}) correspondem aos
anos de 2027 a 2033. A estruturação institucional e social começa primeiro e
se estende por todo o período, porque o modelo de gestão precisa acompanhar os
sistemas até o encerramento. O trabalho técnico-social ocorre entre 2027 e
2030, e a implantação física, entre 2027 e 2033. Os sistemas do Programa
Sanepar Rural são entregues até 2031. Os do IAT e o esgotamento das
comunidades seguem de 2028 a 2033, e o esgotamento da área do IDR-Paraná, de
2029 a 2033.

\begin{figure}[htbp]
  \centering
  \caption{Cronograma de execução do Subcomponente 3.2.b}
  \label{fig:cronograma}
  % Cronograma (diagrama de Gantt) do Subcomponente 3.2.b
% Colunas: Ano 0 (coluna 1) a Ano 6 (coluna 7), conforme o MOP.
\begin{ganttchart}[
    hgrid,
    vgrid={*1{draw=gray!40, dotted}},
    x unit=1.15cm,
    y unit title=0.55cm,
    y unit chart=0.62cm,
    title/.append style={fill=azulclaro, draw=azulpsh},
    title label font=\scriptsize\bfseries,
    bar/.append style={fill=azulpsh!75, draw=azulpsh},
    bar height=0.55,
    bar label font=\scriptsize,
    bar label node/.append style={align=right, text width=6.2cm},
    group/.append style={fill=azulpsh, draw=azulpsh},
    group label font=\scriptsize\bfseries,
    group label node/.append style={align=right, text width=6.2cm},
    milestone/.append style={fill=orange, draw=orange!70!black},
    milestone label font=\scriptsize,
    milestone label node/.append style={align=right, text width=6.2cm}
  ]{1}{7}
  \gantttitle{Ano de execução}{7} \\
  \gantttitle{0}{1} \gantttitle{1}{1} \gantttitle{2}{1} \gantttitle{3}{1}
  \gantttitle{4}{1} \gantttitle{5}{1} \gantttitle{6}{1} \\

  \ganttgroup{I. Estruturação institucional e social}{1}{7} \\
  \ganttbar{3.2.2.1 Modelos de gestão}{1}{7} \\
  \ganttbar{3.2.2.2 Trabalho técnico-social}{2}{5} \\

  \ganttgroup{II. Implantação física}{2}{6} \\
  \ganttbar{3.2.2.3 Abastecimento de água}{3}{6} \\
  \ganttbar{3.2.2.4 Esgotamento -- comunidades}{3}{6} \\
  \ganttbar{3.2.2.5 Esgotamento -- área IDR-Paraná}{3}{6} \\
  \ganttbar{3.2.2.6 Sanepar Rural}{2}{6} \\

  \ganttgroup{III. Monitoramento e encerramento}{3}{7} \\
  \ganttbar{Supervisão técnica e monitoramento}{3}{7} \\
  \ganttbar{Relatórios finais e encerramento}{6}{7} \\
  \ganttmilestone{Termo de Recebimento Definitivo}{7}
\end{ganttchart}

  \fonte{Adaptado de \citeonline{custo2026} e do POA 2026--2027.}
\end{figure}

O Plano Operativo Anual (POA) detalha, mês a mês, o que acontecerá em 2026 e
2027 (Anexo~\ref{an:cronograma}). O ano de 2026 será dedicado à preparação.
A equipe interna do IAT fará o planejamento integrado, a seleção e a adesão
dos municípios, a assinatura dos Termos de Cooperação, a definição do modelo
de gestão piloto e o estudo de alternativas para o esgotamento. No segundo
semestre de 2026, começarão a ser elaborados os primeiros termos de
referência. Em 2027, as contratações estarão em licitação, e as primeiras
execuções estão previstas para o fim do ano. A mobilização social e os
estudos do poço começarão em outubro, a perfuração em novembro e os projetos
das redes em dezembro de 2027. As demais contratações não iniciam a execução em
2027. As aquisições de materiais para as redes de água e para o esgotamento,
inclusive na frente do IDR-Paraná, e a supervisão das obras estarão em
licitação até dezembro de 2027. O segundo lote de materiais de água e a
infraestrutura do modelo de gestão só terão os TR iniciados no fim de 2027.

\pendencia{Atualizar no MOP todas as referências a ``seis anos'' e aos
``Anos 0 a 6'' para o prazo de sete anos (2027 a 2033).}



% ---------------------------------------------------------------------
% Elementos pós-textuais
% ---------------------------------------------------------------------
\postextual

\bibliography{referencias}

\begin{apendicesenv}
\partapendices
% =====================================================================
%  Apêndice A -- Pendências identificadas na reestruturação
%  A lista é gerada automaticamente a partir dos comandos \pendencia{}.
% =====================================================================
\chapter{Pendências identificadas na reestruturação do MOP}
\label{ap:pendencias}

Durante a conversão do item B.3.2.2 do MOP para a estrutura metodológica,
foram encontradas as inconsistências e lacunas listadas abaixo. Cada item
leva ao ponto do texto em que foi registrado. Depois de resolvidos, os
comandos \verb|\pendencia{...}| devem ser removidos do texto, ou a exibição
pode ser desativada em \texttt{main.tex}.

\ifmostrarpendencias
  \listapendencias
\else
  Não há pendências exibidas nesta versão.
\fi

% =====================================================================
%  Apêndice B -- Quadros detalhados de etapas, produtos e evidências
% =====================================================================
\chapter{Quadros de etapas, produtos e evidências}
\label{ap:quadros}

Este apêndice reúne os quadros detalhados do MOP que orientam a execução e a
verificação das atividades descritas no \autoref{cap:3-2-b}. O
\autoref{qua:etapas-3221} detalha a implantação dos modelos de gestão
(\autoref{subsec:modelos-gestao}). O \autoref{qua:etapas-3223} detalha a
implantação dos sistemas coletivos de abastecimento de água, da perfuração do
poço ao encerramento (Seções~\ref{subsec:captacao} a \ref{subsec:entrega}).

\begin{landscape}
\begin{quadro}[p]
\caption{Etapas de implementação dos modelos de gestão}
\label{qua:etapas-3221}
\scriptsize
\renewcommand{\arraystretch}{1.25}
\begin{tabularx}{\linewidth}{|P{2.6cm}|L|P{2.1cm}|Q{1.3cm}|P{2.5cm}|P{2cm}|P{2.6cm}|Q{1.4cm}|P{1.7cm}|}
\hline
\cab{Etapa} & \cab{Descrição técnica} & \cab{Responsável} & \cab{Prazo} &
\cab{Produto esperado} & \cab{Unidade} & \cab{Evidência} & \cab{Periodic.} &
\cab{Validação} \\ \hline
1. Estudos preliminares &
1.1 Seleção de instrumentos de gestão; 1.2 Definição de parceiros, critérios e
indicadores; 1.3 Meios de comunicação dos modelos; 1.4 Contato inicial com o
município; 1.5 Contato inicial com a comunidade; 1.6 Formalização de pré-adesão &
IAT / IDR / Prefeitura & Ano 0 a 1 & Contratos de pré-adesão &
Número de associados & Contratos assinados & Lote & IAT \\ \hline
2. Implantação &
2.1 Elaboração de documentos (estatuto, manuais, cartilhas etc.);
2.2 Enquadramento jurídico da proposta; 2.3 Adesão familiar;
2.4 Regularização da associação &
IAT / Prefeitura & Ano 1 a 2 & Associações regularizadas &
Número de associações & Registro das associações & Lote & IAT \\ \hline
3. Gestão &
3.1 Manuais de desempenho operacional; 3.2 Investimento para ação em EC;
3.3 Estrutura para monitoramento; 3.4 Treinamento &
IAT / Prefeitura / Associação & Ano 3 a 5 & Manuais, estrutura e treinamento &
Pessoas capacitadas & Lista de presença e certificados & Lote & IAT \\ \hline
4. Supervisão técnica &
4.1 Acompanhamento da implantação &
IAT / Prefeitura / Associação & Ano 2 a 6 & Supervisão e gerenciamento de obras &
Relatório de supervisão & Laudos de vistoria e testes de estanqueidade &
Contínua & IAT \\ \hline
5. Monitoramento e avaliação &
5.1 Avaliação de desempenho do processo de implementação da gestão nas
comunidades &
IAT & Ano 2 a 6 & Relatórios de conformidade técnica e ambiental &
Relatório & Dados no SIGARH / GeoPR & Semestral & UGP / MGAS \\ \hline
6. Relatórios e encerramento &
6.1 Consolidação de resultados, relatórios finais e prestação de contas &
IAT & Ano 5 e 6 & Relatório final de gerenciamento e obras concluídas &
Documento final & Termo de Recebimento Definitivo (TRD) & Única &
UGP / Banco Mundial \\ \hline
\end{tabularx}
\fonte{Adaptado de \citeonline[Quadro 28]{mop2026}.}
\end{quadro}
\end{landscape}

\begin{landscape}
\begin{quadro}[p]
\caption{Etapas de implantação dos sistemas coletivos de abastecimento de água}
\label{qua:etapas-3223}
\scriptsize
\renewcommand{\arraystretch}{1.25}
\begin{tabularx}{\linewidth}{|P{2.6cm}|L|P{2.1cm}|Q{1.3cm}|P{2.8cm}|P{2cm}|P{2.4cm}|Q{1.4cm}|P{1.7cm}|}
\hline
\cab{Etapa} & \cab{Descrição técnica} & \cab{Responsável} & \cab{Prazo} &
\cab{Produto esperado} & \cab{Unidade} & \cab{Evidência} & \cab{Periodic.} &
\cab{Validação} \\ \hline
1. Perfuração de poço &
1.1 Estudo hidrogeológico com locação, projeto construtivo do poço e anuência
prévia; 1.2 Obra de perfuração; 1.3 Teste de vazão; 1.4 Análise
físico-química; 1.5 Outorga; 1.6 Gerenciamento e fiscalização de obras &
IAT / Prefeituras & Ano 0 a 1 &
83 locações e projetos, anuências, testes de vazão, análises físico-químicas e
outorgas & Poços perfurados e outorgados &
Atesto de recebimento e laudos de vistoria & Lote & IAT \\ \hline
2. Instalação do sistema de abastecimento &
2.1 Aquisição de projeto básico e executivo (casa química, reservatório,
eletromecânica e distribuição); 2.2 Aprovação dos projetos; 2.3 Obras de
instalação; 2.4 Projeto \textit{as built}; 2.5 Pré-operação, testes e
capacitação dos operadores locais; 2.6 Gerenciamento e fiscalização das obras &
IAT / Prefeituras & Ano 0 a 1 &
83 projetos aprovados, sistemas de armazenamento e distribuição; XXX
hidrômetros instalados & Hidrômetros instalados &
Atesto de recebimento e laudos de vistoria & Lote & IAT \\ \hline
3. Capacitação dos beneficiários &
3.1 Aquisição de treinamentos sobre operação e manutenção dos sistemas &
IAT / Prefeituras & Ano 4 a 6 & Treinamentos na comunidade ou município &
Pessoas capacitadas & Lista de presença e certificados & Contínua & IAT \\ \hline
4. Supervisão técnica &
4.1 Acompanhamento das instalações e suporte técnico pós-implantação &
IAT / Fiscalização de obras & Ano 2 a 6 & Supervisão e gerenciamento de obras &
Relatório de supervisão & Laudos de vistoria e testes de estanqueidade &
Contínua & IAT \\ \hline
5. Monitoramento e avaliação &
5.1 Avaliação de desempenho dos sistemas, indicadores de impacto e ajustes &
IAT & Ano 2 a 6 & Relatórios de conformidade técnica e ambiental &
Relatório & Dados no SIGARH / GeoPR & Semestral & UGP / MGAS \\ \hline
6. Relatórios e encerramento &
6.1 Consolidação de resultados, relatórios finais e prestação de contas &
IAT & Ano 5 e 6 & Relatório final de gerenciamento e obras concluídas &
Documento final & Termo de Recebimento Definitivo (TRD) & Única &
UGP / Banco Mundial \\ \hline
\end{tabularx}
\fonte{Adaptado de \citeonline[Quadro 29]{mop2026}.}
\end{quadro}
\end{landscape}

\end{apendicesenv}

\begin{anexosenv}
\partanexos
% =====================================================================
%  Anexo A -- Cronograma das atividades em 2026 e 2027 (POA)
%  Fonte: POA_PSH_2026_2027_conferido_EAP_0608_V2.xlsx, aba "Componente 3"
%  (linhas 15 a 29 -- Subcomponente 3.2, exceto a linha 14, que trata da
%  política e do plano estadual de saneamento rural, Subcomponente 3.2.a)
% =====================================================================
\chapter{Cronograma das atividades em 2026 e 2027}
\label{an:cronograma}

Este anexo reproduz, para o Subcomponente 3.2.b, o Plano Operativo Anual
(POA) do PSH-PR para 2026 e 2027, conferido com a Estrutura Analítica do
Projeto (EAP) em 31 de julho de 2026. O \autoref{qua:poa-atividades} lista as
15 atividades previstas, com a atividade correspondente do MOP, o tipo de
contratação, o responsável, o marco de conclusão, as datas previstas e a
situação. A \autoref{fig:gantt-poa} mostra, mês a mês, as fases de cada uma.

Nas contratações financiadas pelo BIRD, o POA adota uma regra única. A barra
de TR corresponde aos três meses de elaboração até o TR ficar pronto para
envio ao Banco Mundial. A barra de licitação vai do TR pronto até o mês
anterior ao início da execução, e a barra de execução começa na data prevista
para o início dos serviços (\autoref{fig:ciclo-contratacao}). As atividades
da equipe interna do IAT aparecem como uma barra única, que indica o período
de trabalho previsto em 2026 e 2027.

\begin{landscape}
\begin{quadro}[p]
\caption{Atividades do Subcomponente 3.2.b no POA 2026--2027}
\label{qua:poa-atividades}
\scriptsize
\renewcommand{\arraystretch}{1.2}
\begin{tabularx}{\linewidth}{|Q{0.5cm}|L|Q{1.6cm}|P{2.0cm}|P{2.1cm}|P{3.0cm}|Q{2.2cm}|Q{1.5cm}|}
\hline
\cab{Nº} & \cab{Atividade (POA)} & \cab{Atividade MOP} & \cab{Contratação} &
\cab{Responsável (participantes)} & \cab{Marco} & \cab{Início--fim previstos} &
\cab{Situação} \\ \hline
\multicolumn{8}{|l|}{\cellcolor{cinzaclaro}\textbf{Preparação institucional}} \\ \hline
1 & Planejamento integrado das ações de saneamento rural para 2027 & 3.2.2.1 e 3.2.2.2 &
  Equipe interna & IAT (IDR) & Planejamento entregue & 05/2026--05/2029 & Em andamento \\ \hline
2 & Seleção dos municípios com base nos critérios do programa & 3.2.2.1 &
  Equipe interna & IAT & Mapa & 05/2026--05/2029 & Em andamento \\ \hline
3 & Articulação e formalização da adesão dos municípios & 3.2.2.1 &
  Equipe interna & IAT & Lista dos municípios & 05/2026--05/2029 & Em andamento \\ \hline
4 & Formalização de Termo de Cooperação junto ao município & 3.2.2.1 &
  Equipe interna & IAT & Termo de Cooperação assinado & 05/2026--05/2029 & Em andamento \\ \hline
5 & Articulação para definição do modelo de gestão junto aos entes envolvidos & 3.2.2.1 &
  Equipe interna & IAT (IDR e Sanepar) & Modelo de gestão piloto definido para implementação &
  04/2030--04/2033 & Em andamento \\ \hline
\multicolumn{8}{|l|}{\cellcolor{cinzaclaro}\textbf{Abastecimento de água}} \\ \hline
6 & Mobilização social, estudo hidrogeológico e anuência prévia, teste de vazão, análise
  físico-química, anteprojeto, levantamento preliminar social e cadastro das famílias (busca ativa) &
  3.2.2.2 e 3.2.2.3 & Regras BIRD & IAT e IDR (IDR) & Contrato assinado & 09/2026--01/2029 & Não iniciado \\ \hline
7 & Aquisição de materiais para perfuração de poços & 3.2.2.3 &
  Regras BIRD & IAT & Contrato assinado & 10/2026--11/2031 & Não iniciado \\ \hline
8 & Contratação de projetos das redes de abastecimento de água e de esgotamento sanitário &
  3.2.2.3 e 3.2.2.4 & Regras BIRD & IAT & Contrato assinado & 11/2026--12/2028 & Não iniciado \\ \hline
9 & Aquisição de 25 unidades de materiais para implantação da rede de abastecimento & 3.2.2.3 &
  Regras BIRD & IAT & Contrato assinado & 12/2026--09/2030 & Não iniciado \\ \hline
10 & Aquisição de 30 unidades de materiais para implantação da rede de abastecimento & 3.2.2.3 &
  Regras BIRD & IAT & Contrato assinado & 12/2027--12/2031 & Não iniciado \\ \hline
11 & Supervisão das obras de implantação de SAA e SESI & 3.2.2.3 e 3.2.2.4 &
  Regras BIRD & IAT & Contrato assinado & 12/2026--05/2032 & Não iniciado \\ \hline
\multicolumn{8}{|l|}{\cellcolor{cinzaclaro}\textbf{Esgotamento sanitário}} \\ \hline
12 & Estudo das melhores alternativas de sistemas de esgotamento sanitário em comunidades rurais &
  3.2.2.4 & Equipe interna & IAT (IDR) & Estudo concluído & 10/2026--10/2029 & Não iniciado \\ \hline
13 & Aquisição de material para implantação do esgotamento sanitário & 3.2.2.4 &
  Regras BIRD & IAT & Contrato assinado & 01/2027--04/2032 & Não iniciado \\ \hline
14 & Aquisição de materiais e implantação dos sistemas de esgotamento em domicílios rurais
  (área de atuação do IDR) & 3.2.2.5 & Regras BIRD & IAT (IDR) & Contrato assinado &
  04/2027--05/2032 & Não iniciado \\ \hline
\multicolumn{8}{|l|}{\cellcolor{cinzaclaro}\textbf{Modelo de gestão}} \\ \hline
15 & Aquisição de infraestrutura para implantação do modelo de gestão & 3.2.2.1 &
  Regras BIRD & IAT & Contrato assinado & 11/2027--12/2031 & Não iniciado \\ \hline
\end{tabularx}
\fonte{Adaptado do POA 2026--2027 do PSH-PR (aba ``Componente 3'', linhas 15 a 29), conferido com a EAP em 31/07/2026.}
\end{quadro}
\end{landscape}

\begin{landscape}
\begin{figure}[p]
  \centering
  \caption{Cronograma mensal das atividades do Subcomponente 3.2.b em 2026 e 2027}
  \label{fig:gantt-poa}
  % Cronograma mensal 2026-2027 do Subcomponente 3.2.b (POA)
% Colunas 1 a 12 = jan a dez/2026; 13 a 24 = jan a dez/2027.
\newcommand{\fTR}[3]{\ganttbar[bar/.append style={fill=corTR}]{#1}{#2}{#3}}
\newcommand{\fLic}[2]{\ganttbar[bar/.append style={fill=corLicitacao}]{}{#1}{#2}}
\newcommand{\fExe}[2]{\ganttbar[bar/.append style={fill=corExecucao}]{}{#1}{#2}}
\newcommand{\fInt}[3]{\ganttbar[bar/.append style={fill=corInterna}]{#1}{#2}{#3}}
\begin{ganttchart}[
    hgrid={*1{draw=gray!30}},
    vgrid={*1{draw=gray!25, dotted}},
    x unit=0.6cm,
    y unit title=0.5cm,
    y unit chart=0.52cm,
    title/.append style={fill=azulclaro, draw=azulpsh},
    title label font=\scriptsize\bfseries,
    bar/.append style={draw=gray!60!black},
    bar height=0.6,
    bar label font=\scriptsize,
    bar label node/.append style={align=right, text width=9.2cm},
    group/.append style={fill=none, draw=none},
    group label font=\scriptsize\bfseries\color{azulpsh},
    group label node/.append style={align=right, text width=9.2cm},
    group right shift=0, group top shift=0.45, group height=0.1,
    group peaks height=0
  ]{1}{24}
  \gantttitle{2026}{12} \gantttitle{2027}{12} \\
  \gantttitle{J}{1} \gantttitle{F}{1} \gantttitle{M}{1} \gantttitle{A}{1}
  \gantttitle{M}{1} \gantttitle{J}{1} \gantttitle{J}{1} \gantttitle{A}{1}
  \gantttitle{S}{1} \gantttitle{O}{1} \gantttitle{N}{1} \gantttitle{D}{1}
  \gantttitle{J}{1} \gantttitle{F}{1} \gantttitle{M}{1} \gantttitle{A}{1}
  \gantttitle{M}{1} \gantttitle{J}{1} \gantttitle{J}{1} \gantttitle{A}{1}
  \gantttitle{S}{1} \gantttitle{O}{1} \gantttitle{N}{1} \gantttitle{D}{1} \\

  \ganttgroup{Preparação institucional}{1}{24} \\
  \fInt{1. Planejamento integrado das ações de 2027}{5}{10} \\
  \fInt{2. Seleção dos municípios (mapa)}{5}{12} \\
  \fInt{3. Articulação e adesão dos municípios}{5}{12} \\
  \fInt{4. Termo de Cooperação com o município}{5}{12} \\
  \fInt{5. Definição do modelo de gestão piloto}{4}{12} \\

  \ganttgroup{Abastecimento de água}{1}{24} \\
  \fTR{6. Mobilização social, estudos do poço e cadastro}{9}{11} \fLic{12}{21} \fExe{22}{24} \\
  \fTR{7. Materiais para perfuração de poços}{10}{12} \fLic{13}{22} \fExe{23}{24} \\
  \fTR{8. Projetos das redes de água e esgoto}{11}{13} \fLic{14}{23} \fExe{24}{24} \\
  \fTR{9. Materiais para rede de abastecimento (25 un.)}{12}{14} \fLic{15}{24} \\
  \fTR{10. Materiais para rede de abastecimento (30 un.)}{24}{24} \\
  \fTR{11. Supervisão das obras de água e esgoto}{12}{14} \fLic{15}{24} \\

  \ganttgroup{Esgotamento sanitário}{1}{24} \\
  \fInt{12. Estudo de alternativas de esgotamento}{10}{12} \\
  \fTR{13. Material para esgotamento nas comunidades}{13}{15} \fLic{16}{24} \\
  \fTR{14. Materiais e implantação -- área do IDR-Paraná}{16}{18} \fLic{19}{24} \\

  \ganttgroup{Modelo de gestão}{1}{24} \\
  \fTR{15. Infraestrutura para o modelo de gestão}{23}{24}
\end{ganttchart}

\vspace{0.3cm}
{\scriptsize
\tikz\fill[corInterna, draw=gray!60!black] (0,0) rectangle (0.5,0.25); Equipe interna do IAT \quad
\tikz\fill[corTR, draw=gray!60!black] (0,0) rectangle (0.5,0.25); Elaboração do TR \quad
\tikz\fill[corLicitacao, draw=gray!60!black] (0,0) rectangle (0.5,0.25); Processo licitatório \quad
\tikz\fill[corExecucao, draw=gray!60!black] (0,0) rectangle (0.5,0.25); Execução
}

  \fonte{Adaptado do POA 2026--2027 do PSH-PR (aba ``Componente 3'').}
\end{figure}
\end{landscape}

\pendencia{No POA, a linha da definição do modelo de gestão (nº~5) tem início e
fim previstos em 04/2030 e 04/2033, mas o cronograma marca o período de abril
a dezembro de 2026. As atividades internas (nº~1 a~4 e~12) têm fim previsto
em 2029, embora suas barras terminem em 2026. Conferir as datas.}

\pendencia{A contratação dos projetos das redes (nº~8) não tem correspondência
na EAP, segundo a própria aba de conferência do POA, e tem meses de janeiro a
outubro de 2026 preenchidos sem identificação de fase. O Programa Sanepar
Rural (\atividade{3.2.2.6}) não aparece no POA do Subcomponente 3.2.}

\end{anexosenv}

\end{document}
